The following result is one of the most useful in the paper. Intuitively, it
says that in projective mini-schemes, a counterfeiter that copies a banknote
with \emph{any} non-negligible fidelity can be \textquotedblleft
amplified\textquotedblright\ to a counterfeiter that copies the banknote
almost \emph{perfectly}---or conversely, that to rule out the former sort of
counterfeiter, it suffices to rule out the latter. The proof makes essential
use of the amplitude amplification results from \expref{Section}{SEARCH}.

\begin{theorem}
[Amplification of Counterfeiters]\label{miniamp}Let $\mathcal{M}=\left(
\mathsf{Bank},\mathsf{Ver}\right)  $ be a projective mini-scheme, and let
$\$=\left(  s,\rho\right)  $ be a valid banknote in $\mathcal{M}$. Suppose
there exists a counterfeiter $C$ that copies $\$$ with probability
$\varepsilon>0$:\ that is,%
\[
\Pr\left[  \mathsf{Ver}_{2}\left(  s,C\left(  \$\right)  \right)  \text{
accepts}\right]  \geq\varepsilon.
\]
Then for all $\delta>0$, there is also a modified counterfeiter $C^{\prime}$
(depending only on $\varepsilon$ and $\delta$, not $\$$), which makes%
\[
O\left(  \frac{\log1/\delta}{\sqrt{\varepsilon}\left(  \sqrt{\varepsilon
}+\delta^{2}\right)  }\right)
\]
queries to $C$, $C^{-1}$, and $\mathsf{Ver}$ and which satisfies%
\[
\Pr\left[  \mathsf{Ver}_{2}\left(  s,C^{\prime}\left(  \$\right)  \right)
\text{ accepts}\right]  \geq1-\delta.
\]

\end{theorem}

\begin{proof}
Write $\$$ as a mixture of pure states:%
\[
\$=\sum p_{i}\left\vert \psi_{i}\right\rangle \left\langle \psi_{i}\right\vert
.
\]
By linearity, clearly it suffices to show that%
\[
\Pr\left[  \mathsf{Ver}_{2}\left(  s,C^{\prime}\left(  \left\vert \psi
_{i}\right\rangle \right)  \right)  \text{ accepts}\right]  \geq1-\delta
\]
for all $i$ such that $p_{i}>0$. We focus on $\left\vert \psi\right\rangle
:=\left\vert \psi_{1}\right\rangle $ without loss of generality.

By assumption, there exists a subspace $S$ such that%
\[
\Pr\left[  \mathsf{Ver}\left(  \rho\right)  \text{ accepts}\right]  =F\left(
\rho,S\right)  ^{2}%
\]
for all $\rho$. Then $F\left(  \$,S\right)  =F\left(  \left\vert
\psi\right\rangle ,S\right)  =1$.

Now, just as $\mathsf{Ver}$ is simply a projector onto $S$, so $\mathsf{Ver}%
_{2}$ is a projector onto $S^{\otimes2}$. Thus%
\[
F\left(  C\left(  \left\vert \psi\right\rangle \right)  ,S^{\otimes2}\right)
\geq\sqrt{\varepsilon}.
\]
So consider performing a fixed-point Grover search, with $C\left(  \left\vert
\psi\right\rangle \right)  $ as the initial state and $S^{\otimes2}$ as the
target subspace. By \expref{Lemma}{fixedpoint}, this will produce a state $\rho
$ such that $F\left(  \rho,S^{\otimes2}\right)  \geq1-\delta$ using
$O\left(({1}/{\varepsilon}) \log ({1}/{\delta})\right)  $ Grover
iterations. Each iteration requires a reflection about $C\left(  \left\vert
\psi\right\rangle \right)  $ and a reflection about $S^{\otimes2}$, which can
be implemented using $O\left(  1\right)  $ queries to $C,C^{-1}$ and
$\mathsf{Ver}$ respectively. Therefore the number of queries to $C,C^{-1}$
and $\mathsf{Ver}$ is $O\left(({1}/{\varepsilon}) \log ({1}/{\delta
})\right)$ as well.

If $\delta$ is large compared to $\varepsilon$, then we can instead use
\expref{Theorem}{combined}, which produces a state $\rho$ such that $F\left(
\rho,S^{\otimes2}\right)  \geq1-\delta$ using 
\[
O\left(  \frac{1}
{\sqrt{\varepsilon}\delta^{2}}\log\frac{1}{\delta}\right)
\]
iterations.
Taking the minimum of the two bounds gives us the claimed bound on query complexity.
\end{proof}
